\documentclass[11pt,a4paper]{article}

\usepackage[utf8]{inputenc}
\usepackage[T1]{fontenc}
\usepackage{lmodern}
\usepackage[english]{babel}
\usepackage{geometry}
\usepackage{graphicx}
\usepackage{float}
\usepackage{caption}
\usepackage{amsmath}
\usepackage{amssymb}
\usepackage{booktabs}
\usepackage{hyperref}
\usepackage{subcaption}
\usepackage{xcolor}

\geometry{margin=1in}
\raggedbottom
\emergencystretch=3em

% Figures live in the repo's runs/ directory; this file sits in report/7/,
% so resolve \includegraphics paths relative to the repository root.
\graphicspath{{../../}}

% Graceful fallback for figures produced by runs that may not be pulled yet.
\newcommand{\figorbox}[2]{%
  \IfFileExists{../../#1}{\includegraphics[width=#2\textwidth]{#1}}{%
    \fbox{\parbox[c][3.5cm][c]{#2\textwidth}{\centering\itshape
    \textcolor{gray}{[figure not yet generated --- run the eval and pull it]}}}}}

\title{Empirical Analysis of Transformers on Graph Connectivity \\
       \large Part VII: Opening the Trunk --- Why Does an Unbalanced Two-Chain Generalise?}
\author{Edoardo Paccagnella}
\date{\today}

\begin{document}

\maketitle

% =====================================================================
% DRAFT / SKELETON. No data collected yet. This file lays out the framing and
% the planned experiments for Report VII. Each results subsection is a STUB:
% the % comments record (a) the question, (b) which hypothesis it adjudicates,
% (c) the exact checkpoints/test distribution to use, (d) the caption fields
% required by rule 21. Replace the italic "[To run --- data pending]" lines
% with real tables/figures once the runs are in.
%
% Design rules carried over from istruzioni.md (esp. errors 21, 55-60):
%  - every caption states MODEL + TRAINING SET + TEST SET + METRIC, no analysis;
%  - no code/file/path/flag names in the rendered text;
%  - no 'advisor' attributions, no chronological "we re-ran" story;
%  - per-seed figures/tables wherever an effect could be seed-bimodal, not
%    mean+band;
%  - state hypotheses up front and report which survive, honestly.
% =====================================================================

\begin{abstract}
\noindent Report~VI (Thread~B) found a puzzle it could not explain: the same
trained $L{=}2$ transformer, same weights, same canvas size $n{=}40$, solves a
two-chain graph split into paths of length $a$ and $n-a$ perfectly when $a$ is
small (e.g.\ $4{+}36$, correct out to the far end of the long path) but fails
completely once the split is close to balanced (e.g.\ $17{+}23$, $20{+}20$ ---
correct only up to the architectural distance wall $3^L{=}9$, then exactly zero).
Every report so far has read the model \emph{behaviourally}: from its outputs, not
from what it computes. This report opens the trunk --- the attention weights, the
read-out matrix, the per-layer embeddings --- and asks what is mechanistically
different between a split the model solves and one it does not.
\end{abstract}

\section{Introduction}
\label{sec:intro}

\subsection{What Reports I--VI established}
\label{sec:recap}

We treat the following as established (documented, multi-seed, in Reports~I--VI);
the new investigation must stay consistent with them.

\begin{itemize}\setlength\itemsep{0.2em}
  \item \textbf{There is a sharp, architectural distance wall at $3^L$} ($=9$ at
        $L=2$), which moves with depth and appears even on graph families seen
        $\sim10^8$ times in training (Reports~III--IV).
  \item \textbf{The model computes distance-bounded matrix powering, not a bounded
        traversal}: a DFS/BFS reading is rejected on constructed and (partly) on
        natural graphs; a soft, learned ``node budget'' on dense structure is a
        removable data prior, not a second capacity (Report~V).
  \item \textbf{Parallel routes and repeated local structure both help, within
        capacity, in ways that are data-shaped rather than architecturally free}:
        more disjoint routes between two endpoints rescue a connection the model
        could not otherwise confirm; a learned local rule (a bridge between two
        cliques) composes only weakly along a chain (Report~VI).
  \item \textbf{The unexplained result this report targets}: sweeping the split
        $(a, n-a)$ of a two-chain graph at fixed $n=40$, the same checkpoint goes
        from solving \emph{every} pair in the long path (out to $\sim38$ hops) to
        solving \emph{none} beyond the $9$-hop wall, as $a$ crosses a threshold
        close to the capacity itself (Report~VI, Thread~B). No account of the
        mechanism was given.
  \item \textbf{Every report through Report~VI is behavioural.} Model outputs are
        compared against oracles, ablated data distributions, or read-out variants
        --- never against the attention weights or the read-out matrix themselves.
        \texttt{model.py} exposes tools for exactly that (real per-layer attention
        weights, per-layer embeddings) that no report has yet used.
\end{itemize}

\subsection{The new question: what is mechanistically different?}
\label{sec:question}

Report~VI's own closing paragraph names this split puzzle as the one open thread it
could reproduce but not explain. This report is the direct follow-up: instead of
another behavioural sweep, we open the trained model at the level of its attention
scores and its read-out weights, on the exact same checkpoints, and ask what
changes between a split the model solves and one it does not.

\subsection{Three candidate hypotheses}
\label{sec:hypotheses}

We state candidate explanations up front, so each experiment is a test that could
come out either way --- not a search for evidence of an assumed answer.

\begin{enumerate}\setlength\itemsep{0.15em}
  \item \textbf{Genuine propagation.} The model actually traces distance along the
        path; behaviour should be robust to relabelling nodes, to where the short
        component sits, and to the internal edge order of the long path.
  \item \textbf{Short-component-resolved default.} The model fully resolves
        whichever component is small enough to fit inside its capacity
        ($\mathrm{short\_len}\lesssim9$) and defaults every other pair to
        ``connected'' --- a shortcut that is correct whenever exactly one
        component is small, and has nothing to fall back on once \emph{both}
        components are too large to resolve.
  \item \textbf{Label/order shortcut.} The model exploits some regularity in how
        node indices map to structure, rather than the abstract graph. (The
        existing evaluation already relabels every test graph independently at
        random before scoring it, which substantially weakens the crude form of
        this hypothesis before we even start; we still check it directly against
        the read-in/read-out weights.)
\end{enumerate}

\section{Setup}
\label{sec:setup}

\paragraph{Model.} The same base model as Reports~III--VI: the RoBERTa-faithful
standard transformer with the \emph{linear} read-out, $L=2$, single head,
$d_{\mathrm{model}}=512$, normalised-ReLU attention $\alpha=(1/n)\,\mathrm{ReLU}
(QK^\top/\sqrt{d_h})$, trained online at $n=40$ on a single clean distribution of
disjoint paths (unions of $1$--$4$ paths covering all $n$ nodes), four seeds. This
is the exact checkpoint set behind Report~VI's Thread~B.

\paragraph{The forward pass, in full.} Because \S\ref{sec:res-attention}--\S\ref{sec:res-heatmaps}
open the trunk directly, we fix here, once, exactly what the model computes at inference, so that
every quantity used later --- the attention score, the normalised-ReLU weights, the message
contribution, the node-to-node contribution --- has a precise definition. The input is the
self-loop-augmented adjacency $A'=A+I\in\{0,1\}^{n\times n}$. The read-in, once,
\[ H^{(0)} = \mathrm{LayerNorm}\big(A'\,W_{\mathrm{in}} + \mathbf 1\, b_{\mathrm{in}}^\top\big), \]
starts row $i$, before the LayerNorm, as simply the sum of one learned vector per neighbour of
$i$ (including $i$ itself, via the self-loop) plus a bias --- no attention yet. Each of the two
transformer blocks (layer $\ell=1,2$, post-LayerNorm) then maps $H^{(\ell-1)}$ to $H^{(\ell)}$ by
\begin{align*}
Q_\ell &= H^{(\ell-1)}W_Q^{(\ell)}, \quad K_\ell = H^{(\ell-1)}W_K^{(\ell)}, \quad
V_\ell = H^{(\ell-1)}W_V^{(\ell)}\\[2pt]
S_\ell &= Q_\ell K_\ell^\top/\sqrt{d_h}\\[2pt]
\alpha_\ell &= \tfrac1n\,\mathrm{ReLU}(S_\ell)\\[2pt]
\mathrm{AttnOut}_\ell &= (\alpha_\ell V_\ell)\,W_O^{(\ell)}\\[2pt]
\widehat H_\ell &= \mathrm{LayerNorm}\big(H^{(\ell-1)}+\mathrm{AttnOut}_\ell\big)\\[2pt]
H^{(\ell)} &= \mathrm{LayerNorm}\big(\widehat H_\ell + \mathrm{FFN}^{(\ell)}(\widehat H_\ell)\big)
\end{align*}
Here $d_h=d_{\mathrm{model}}$ (single head); $S_\ell$ is the raw attention score; $\alpha_\ell$
is the normalised-ReLU attention --- \emph{not} a softmax, it does not sum to $1$ per row;
$\alpha_\ell$ weights the \emph{values} $V_\ell$, and $W_O^{(\ell)}$ is the output projection
applied afterwards; $\widehat H_\ell$ adds that attention output back onto the residual stream;
$H^{(\ell)}$ then adds the feed-forward step $\mathrm{FFN}^{(\ell)}$ on top, again as a residual.
Layer $2$ repeats this starting from $H^{(1)}$, with its \emph{own} $W_Q,W_K,W_V,W_O$ and its own
feed-forward weights --- it does not reuse $\alpha_1$ in any way; the two layers are two separate
passes of this same recipe, run one after the other. The read-out (\S\ref{sec:setup} below) is
applied once, at the very end, to $H^{(2)}$.

\paragraph{The same computation, one number at a time.} Every equation above is a matrix
equation; every one of them is also just an ordinary sum over the $d:=d_{\mathrm{model}}=512$
hidden features, and it is worth writing out explicitly, matching the actual implementation
entry by entry rather than only in matrix form. Two building blocks recur throughout and are
defined once, here. A linear layer with weight matrix $W$ (indexed input-feature, output-feature)
and bias $b$ maps a row $x$ to
\[ y_r \;=\; \sum_{s} x_s\,W_{s,r} \;+\; b_r; \]
every one of $W_{\mathrm{in}}, W_Q^{(\ell)}, W_K^{(\ell)}, W_V^{(\ell)}, W_O^{(\ell)}, W_1^{(\ell)},
W_2^{(\ell)}$ in this report is one such matrix, each with its own bias --- all dropped from the
matrix equations above only to keep them short, never dropped from the actual computation.
\textbf{LayerNorm} is applied independently to each \emph{node's own} row $x\in\mathbb R^d$, with
its own learned per-feature scale $\gamma$ and shift $\beta$:
\[ \mathrm{LayerNorm}(x)_r \;=\; \gamma_r\,\frac{x_r-\bar x}{\sqrt{\mathrm{Var}(x)+\epsilon}} \;+\; \beta_r,
\qquad \bar x=\frac1d\sum_{s} x_s,\quad \mathrm{Var}(x)=\frac1d\sum_{s}(x_s-\bar x)^2; \]
crucially $\bar x$ and $\mathrm{Var}(x)$ come from that \emph{one} node's own $d$ numbers, never
averaged across nodes, and there are five independent LayerNorms in the trunk (the embedding one,
plus one after each residual, in each of the two layers), each with its own $\gamma,\beta$.
$\mathrm{GELU}(x)=x\,\Phi(x)$ applies elementwise, $\Phi$ the standard normal CDF.

With these two building blocks, every quantity in this report is, entry by entry: the read-in,
\[ u_{i,r} = \sum_{k=1}^n A'_{ik}\,(W_{\mathrm{in}})_{k,r} + (b_{\mathrm{in}})_r,
\qquad h^{(0)}_{i,r} = \mathrm{LayerNorm}(u_i)_r; \]
the query at layer $\ell$ (identically for $k_{i,r}$ and $v_{i,r}$, with their own weights),
\[ q_{i,r} = \sum_{s=1}^{d} h^{(\ell-1)}_{i,s}\,(W_Q^{(\ell)})_{s,r} + (b_Q^{(\ell)})_r; \]
the attention score and weight for the \emph{ordered pair} $(i,j)$ --- a single number, not a
matrix, one for every pair of nodes,
\[ S_{\ell,ij} = \frac{1}{\sqrt{d_h}}\sum_{r=1}^{d_h} q_{i,r}\,k_{j,r},
\qquad \alpha_{\ell,ij} = \frac1n\max(0,\,S_{\ell,ij}); \]
the attention output --- first mix \emph{values} across every other node $j$, feature by feature,
then project the mixed result,
\[ m_{\ell,i,s} = \sum_{j=1}^n \alpha_{\ell,ij}\,v_{j,s},
\qquad \mathrm{AttnOut}_{\ell,i,r} = \sum_{s=1}^{d} m_{\ell,i,s}\,(W_O^{(\ell)})_{s,r} + (b_O^{(\ell)})_r; \]
the first residual and LayerNorm,
\[ \widehat h_{\ell,i,r} = \mathrm{LayerNorm}\big(h^{(\ell-1)}_{i,\cdot} + \mathrm{AttnOut}_{\ell,i,\cdot}\big)_r; \]
the feed-forward step (hidden width $d_{\mathrm{ff}}=2048$, GELU in between the two linear maps),
\[ u^{\mathrm{ff}}_{\ell,i,m} = \sum_{s=1}^{d} \widehat h_{\ell,i,s}\,(W_1^{(\ell)})_{s,m} + (b_1^{(\ell)})_m,
\qquad g_{\ell,i,m} = \mathrm{GELU}\big(u^{\mathrm{ff}}_{\ell,i,m}\big), \]
\[ \mathrm{FFN}^{(\ell)}(\widehat h_{\ell,i})_r = \sum_{m=1}^{d_{\mathrm{ff}}} g_{\ell,i,m}\,(W_2^{(\ell)})_{m,r} + (b_2^{(\ell)})_r; \]
and the second residual and LayerNorm, closing the layer,
\[ h^{(\ell)}_{i,r} = \mathrm{LayerNorm}\big(\widehat h_{\ell,i,\cdot} + \mathrm{FFN}^{(\ell)}(\widehat h_{\ell,i})\big)_r. \]
Every sum above runs over an ordinary feature index ($r,s,m$); the \emph{only} place a sum ever
runs over all $n$ nodes at once, rather than over $d$ or $d_{\mathrm{ff}}$ features, is the
value-mixing step $m_{\ell,i,s}=\sum_j\alpha_{\ell,ij}v_{j,s}$ --- the one place information
actually moves between nodes. The same pattern gives the read-out logit at the level of a single
pair (\S\ref{sec:setup} below): $z_{ij} = \sum_{r=1}^{d} h^{(2)}_{i,r}\,(W_{\mathrm{out}})_{r,j} + b_j$.

\paragraph{Message contribution and node-to-node contribution: what they measure, and how they
are computed.} Two derived quantities recur in \S\ref{sec:res-attention} and
\S\ref{sec:res-heatmaps}. The \textbf{message contribution} $\alpha_{ij}\|W_O v_j\|$ uses the
actual value vector $v_j$ and output projection $W_O$ of one layer to size the additive message
node $j$ sends into node $i$'s residual stream \emph{at that layer} --- it is exact for that one
layer's attention branch, but does not include the residual stream itself or the feed-forward
step, both of which also update $h_i$. The \textbf{node-to-node contribution}
\[ C_{ik} \;=\; \left\|\frac{\partial h_i^{(2)}}{\partial h_k^{(0)}}\right\|_F \]
answers the question directly, across \emph{both} layers at once: how much does node $i$'s final
embedding move if node $k$'s read-in embedding moves, holding everything else fixed? We compute
it by literal automatic differentiation through the exact forward pass above --- the real
$V,W_O$, the real residual identity, the real feed-forward step, and LayerNorm's true local
derivative are all part of the computation, none of them dropped or approximated; the only sense
in which this is ``local'' is the ordinary one for any derivative, that it is evaluated at one
specific graph's actual activations.

\paragraph{Why the read-out shape matters.} With the linear read-out, the logit for
pair $(i,j)$ is $z_{ij} = h_i^\top w_j + b_j$, where $w_j$ is the $j$-th row of the
read-out matrix $W_{\mathrm{out}}$ --- a fixed, learned vector attached to the
\emph{target} node $j$, not derived from $j$'s own embedding. Node $i$'s embedding
$h_i$ alone must therefore carry enough information to answer, for \emph{every}
other node $j$, ``are we connected'' --- all the structure the model has learned
about node $j$ lives in $w_j$, fixed across every graph. This is why the read-out
matrix itself, not only the attention weights, is a first-class object of study
here.

\paragraph{Test graphs.} Two-chain graphs partitioning all $n=40$ nodes into two
disjoint paths of size $a$ and $n-a$, swept over the \emph{full} range $a=1,\dots,20$
(Report~VI reported only a handful of points). A new three-component construction
--- one small, resolvable path plus two large, individually-unresolvable paths that
are \emph{not} connected to each other --- is used as a falsification probe for
hypothesis~2 (\S\ref{sec:res-threeway}).

\section{Plan of experiments}
\label{sec:plan}

\paragraph{Behavioural groundwork (\S\ref{sec:res-sweep}--\ref{sec:res-relabel}).} A
dense split sweep and an explicit random-relabelling check, to fix the shape of the
puzzle precisely and close off the crude label-shortcut hypothesis before opening
the trunk.

\paragraph{Read-out mechanism (\S\ref{sec:res-readout}--\ref{sec:res-wout}).} The
$h_i^\top w_j$ decomposition and the geometry of $W_{\mathrm{out}}$'s rows ---
does the read-out encode component membership as a blob, or a graded function of
distance?

\paragraph{Attention mechanism (\S\ref{sec:res-attention}).} Real per-layer
attention weights, the exact node-to-node contribution, and the message
contribution each node actually sends --- does attention itself reach
differently in a solved split than in a failed one?

\paragraph{Falsification test (\S\ref{sec:res-threeway}).} The three-component
construction: if the model is really defaulting ``not in the small component'' to
connected, it should wrongly merge two large components that are genuinely
disconnected from each other.

\paragraph{A second training distribution, and a third canvas size.} Every one of the four
threads above is repeated, at the same $n=40$ canvas, on a model trained purely on sparse
Erd\H{o}s--R\'enyi graphs --- which never contains a disconnected or multi-path graph in
training --- immediately after its disjoint-paths counterpart, so the two conditions can be
read side by side. \S\ref{sec:res-n64} then moves to a larger, sparser canvas ($n=64$);
\S\ref{sec:res-n20} closes the sweep at a smaller one ($n=20$, where the architectural
capacity can no longer be exceeded by any split), on both training distributions.

\paragraph{Why the read-out itself is also worth varying.} Every result up to this point uses
the linear read-out, $z_{ij}=h_i^\top w_j+b_j$ (\S\ref{sec:setup}). This has a specific
architectural consequence worth stating plainly: $w_j$ is a single vector, fixed once training
ends, that must work as a classifier for node $j$ against \emph{every} possible query
embedding $h_i$ --- it is not itself a function of $j$'s own embedding, only of $j$'s identity
as a target. A finding built on this read-out could, in principle, be a property of that
specific asymmetric construction rather than of how the trunk represents connectivity in
general. The similarity read-out already used for other purposes in Report~IV,
$z_{ij}=\mathrm{scale}\cdot\cos(h_i,h_j)+\mathrm{bias}$, is a genuine function of \emph{both}
endpoints' embeddings and carries no such fixed per-target vector, so repeating this report's
full battery on it is a direct test of whether the two-component completion signal is an
artefact of the linear read-out's shape or a property of the trunk that survives a
structurally different, more symmetric read-out. \S\ref{sec:res-similarity} reports the
result, on both training distributions, at the same $n=40$ canvas.

\section{Results}
\label{sec:results}

\subsection{The split sweep, in full}
\label{sec:res-sweep}

\paragraph{Question.} Report~VI reported exact match at seven splits. Sweeping
every $a=1,\dots,20$ locates the break precisely and lets us compare the
model's aggregate behaviour against a genuine distance-bounded oracle.

\paragraph{Result: the break sits one node above the capacity, and beyond it the
model's aggregate behaviour matches pure distance-bounded matrix powering almost
exactly.} Figure~\ref{fig:r7sweep} (top) confirms Report~VI's shape at full
resolution: exact match is $\approx1.0$ for $a\le7$ (with ordinary sampling
noise at $a=3,6$: $0.80$--$0.87$, consistent across all four seeds, not a
single bad run), collapses to exactly $0$ from $a=8$ on, and the long-component
reach follows the same break with a soft floor-plus-recovery from $a\approx11$
to $a=20$ --- the familiar shape of every beyond-capacity curve in this project.
The new observation is the dashed and dotted curves: the model's overall
predicted-positive rate (fraction of all pairs called ``connected'') tracks the
\emph{true} positive rate closely while $a\le7$, then drops sharply and, from
$a\approx11$ on, sits almost exactly on top of the positive rate a pure
distance-$\le9$ oracle would produce on the same graphs ($0.36$ measured vs.\
$0.35$ from the oracle, flat in $a$ because a long path's $9$-hop neighbourhood
size does not depend on how far away the path's other end is). \textbf{Once
both components are individually beyond capacity, the model's aggregate
behaviour is indistinguishable from plain distance-bounded matrix powering}; the
extra ``complete the whole long component'' behaviour is confined to the regime
where one side is small.

\begin{table}[H]\centering\small
\begin{tabular}{lccccc}
\toprule
split $(a,\,40{-}a)$ & exact & reach (long) & reach (short) & cut & pred.\ positive rate\\
\midrule
$(1,39)$  & 1.000 & 1.000 & --    & 1.000 & 0.951\\
$(4,36)$  & 1.000 & 1.000 & 1.000 & 1.000 & 0.820\\
$(7,33)$  & 0.779 & 1.000 & 1.000 & 0.998 & 0.712\\
$(8,32)$  & 0.000 & 0.992 & 1.000 & 0.925 & 0.699\\
$(9,31)$  & 0.000 & 0.873 & 1.000 & 0.938 & 0.599\\
$(10,30)$ & 0.000 & 0.648 & 1.000 & 0.999 & 0.434\\
$(14,26)$ & 0.000 & 0.581 & 0.889 & 1.000 & 0.362\\
$(17,23)$ & 0.000 & 0.638 & 0.792 & 1.000 & 0.362\\
$(20,20)$ & 0.000 & 0.709 & 0.709 & 1.000 & 0.362\\
\bottomrule\end{tabular}
\caption{\emph{Model:} base standard transformer with the linear read-out
($L=2$, single head, $d_{\mathrm{model}}=512$), trained on a single clean
distribution of disjoint paths, four seeds, $n=40$ (the checkpoints behind
Report~VI's Thread~B). \emph{Test set:} two-chain graphs split $(a,\,40{-}a)$,
$200$ freshly sampled, independently node-relabelled graphs per split per seed
($800$ total per split). \emph{Metrics} (mean over seeds): whole-graph exact
match; pairwise reach inside the long/short path (target connected; ``--'' at
$a=1$, no internal short-path pairs); pairwise accuracy across the two paths
(target disconnected); the model's overall fraction of pairs predicted
connected. Full sweep $a=1,\dots,20$ in Figure~\ref{fig:r7sweep}.}
\label{tab:r7sweep}
\end{table}

\begin{figure}[H]
    \centering
    \figorbox{runs/report7/report7_figs/r7_sweep_and_logit.png}{0.92}
    \caption{\emph{Top: Model/training/test as in Table~\ref{tab:r7sweep}}, full
    split sweep $a=1,\dots,20$; the dashed line is the model's overall predicted-
    positive rate, the dotted line the positive rate a pure distance-$\le9$
    oracle would give on the identical graphs (both computed in closed form from
    the two path lengths, independent of any model). \emph{Bottom: same model
    and test graphs}, the mean raw read-out logit $h_i^\top w_j$ (before adding
    the target's bias $b_j$; \S\ref{sec:setup}) for pairs within the long path,
    within the short path, and across the two paths, mean over seeds. Red dashed
    line: capacity $3^L=9$.}
    \label{fig:r7sweep}
\end{figure}

\paragraph{The collapse is not a discontinuity in the mechanism.} Figure~\ref{fig:r7sweep}
(bottom) is the reason the exact-match curve looks like a cliff: the raw,
pre-bias logit $h_i^\top w_j$ for a pair inside the long path is not a switch,
it is a smoothly decaying quantity. It sits saturated and strongly positive for
$a\le7$ ($\approx17$--$36$), falls through zero across $a=8$--$11$ (exactly
where accuracy collapses), and settles into a small, slowly \emph{recovering}
positive value from $a\approx11$ to $a=20$ --- the same soft floor-plus-recovery
shape as the accuracy curve above it, now visible in the model's raw confidence
rather than only in its thresholded predictions. The cut logit (across the two
paths, target disconnected) tells a different story: it starts very negative
($\approx-21$ at $a=1$) and its magnitude shrinks steadily as the split
balances, but it \emph{never crosses zero} at any split tested --- the
``different component'' decision degrades in confidence but not in sign,
consistent with cut accuracy staying at $\approx1.0$ across the whole sweep in
Table~\ref{tab:r7sweep}.

\subsection{The same sweep on an ER-trained model: the failure mode mirrors, it does not repeat}
\label{sec:res-sweep-er}

\paragraph{Question.} Every result above is on a model trained on disjoint paths. Does a
model that has never seen a two-chain, or any multi-component graph, in training show the
same completion signal on the identical test sweep --- or does it fail a different way?

\paragraph{Result: reach barely degrades at any split, but cut collapses steadily as the
split balances --- the mirror image of the disjoint-paths result.} Table~\ref{tab:r7sweeper}
gives the same sweep on four checkpoints trained purely on sparse Erd\H{o}s--R\'enyi graphs
($n=40$, never a two-chain or any disconnected graph in training). The shape is entirely
different from Table~\ref{tab:r7sweep}. Reach in the long component never really drops: it
sits at $0.90$--$0.98$ across the \emph{whole} sweep, with no floor-and-recovery dip in the
middle --- the model confirms far-apart pairs about as well at a balanced split as at an
unbalanced one. What degrades instead is cut: from $1.000$ at $a=1$ it falls steadily to
$0.15$ by $a=20$, the opposite of Table~\ref{tab:r7sweep}, where cut stayed at
$\approx1.0$ throughout and reach was the one that collapsed. Exact match also breaks earlier
and less cleanly than on the disjoint-paths model (already $0$ by $a=7$, where the
disjoint-paths model was still $0.78$), because even the unbalanced splits solved here are
solved by high reach \emph{and} a still-imperfect cut, not by a clean short-component
resolution. The raw read-out logit (not tabulated) confirms the mechanism directly: the
within-long logit stays modest throughout ($3.8$--$6.5$, never saturating the way the
disjoint-paths model's does at $17$--$36$), while the cut logit --- unlike the disjoint-paths
model, where it never changed sign --- starts very negative ($-25.4$ at $a=1$) and
\emph{crosses zero} around $a\approx14$, turning mildly positive by $a=17$--$20$: the model
does not just lose confidence that a cross-pair is disconnected, for many such pairs it comes
to believe the opposite.

\begin{table}[H]\centering\small
\begin{tabular}{lccccc}
\toprule
split $(a,\,40{-}a)$ & exact & reach (long) & reach (short) & cut & pred.\ positive rate\\
\midrule
$(1,39)$  & 0.108 & 0.964 & --    & 1.000 & 0.918\\
$(4,36)$  & 0.719 & 0.999 & 1.000 & 0.999 & 0.820\\
$(7,33)$  & 0.000 & 0.978 & 1.000 & 0.533 & 0.831\\
$(8,32)$  & 0.000 & 0.947 & 1.000 & 0.556 & 0.789\\
$(9,31)$  & 0.000 & 0.920 & 1.000 & 0.575 & 0.753\\
$(10,30)$ & 0.000 & 0.899 & 1.000 & 0.523 & 0.749\\
$(14,26)$ & 0.000 & 0.924 & 0.993 & 0.292 & 0.835\\
$(17,23)$ & 0.000 & 0.944 & 0.977 & 0.192 & 0.885\\
$(20,20)$ & 0.000 & 0.961 & 0.958 & 0.149 & 0.906\\
\bottomrule\end{tabular}
\caption{\emph{Model:} base standard transformer with the linear read-out ($L=2$, single
head, $d_{\mathrm{model}}=512$), trained on pure sparse Erd\H{o}s--R\'enyi
($\mathrm{ER}(40,0.05)$, mean degree $\approx2.0$; never any disconnected or multi-path
graph), four seeds, $n=40$. \emph{Test set/metrics as in Table~\ref{tab:r7sweep}}.}
\label{tab:r7sweeper}
\end{table}

\begin{figure}[H]
    \centering
    \figorbox{runs/report7/report7_figs/r7_sweep_and_logit_n40_er.png}{0.92}
    \caption{\emph{Model/test/metrics as in Figure~\ref{fig:r7sweep}} but trained on pure
    sparse Erd\H{o}s--R\'enyi ($\mathrm{ER}(40,0.05)$, mean degree $\approx2.0$), four seeds,
    $n=40$.}
    \label{fig:r7sweeper}
\end{figure}

\subsection{Relabelling does not explain it}
\label{sec:res-relabel}

Every graph in \S\ref{sec:res-sweep} is already independently relabelled at
random before being scored (a fresh random permutation per graph, unpermuted
before comparing predictions) --- this is the existing methodology, not a new
control. We add one direct comparison: the same split sweep with node labels
\emph{fixed} to a canonical order (the short path always occupying the first
$a$ indices) instead of randomised. Table~\ref{tab:r7relabel} shows no
meaningful difference at any split. Hypothesis~3 in its crude form --- the model
exploiting a fixed mapping from node index to structure --- is not what is
happening here.

\begin{table}[H]\centering\small
\begin{tabular}{lcc}
\toprule
split $(a,\,40{-}a)$ & exact, random relabelling & exact, fixed labelling\\
\midrule
$(1,39)$  & 1.000 & 1.000\\
$(4,36)$  & 1.000 & 1.000\\
$(7,33)$  & 0.779 & 0.750\\
$(8,32)$  & 0.000 & 0.000\\
$(10,30)$ & 0.000 & 0.000\\
$(17,23)$ & 0.000 & 0.000\\
$(20,20)$ & 0.000 & 0.000\\
\bottomrule\end{tabular}
\caption{\emph{Model/training as in Table~\ref{tab:r7sweep}}. \emph{Test set:}
two-chain graphs at the split shown, $200$ graphs per seed, either
independently relabelled at random per graph (the condition behind every other
result in this report) or with node labels fixed to a canonical order (short
path first). \emph{Metric:} whole-graph exact match, mean over four seeds.}
\label{tab:r7relabel}
\end{table}

\subsection{What the read-out logit actually responds to}
\label{sec:res-readout}

\S\ref{sec:res-sweep} already showed the headline pattern (Figure~\ref{fig:r7sweep},
bottom): the within-long logit is large and saturated while the short
component is resolvable, collapses through the transition, and settles at a
small, recovering value that mirrors distance-bounded matrix powering. The
within-short logit is the mirror image by construction (symmetric around
$a=20$, where the two logits meet exactly). One additional check confirms the
decomposition is measuring the right quantity: aggregating $h_i^\top w_j+b_j>0$
by pair type reproduces the behavioural reach numbers of Table~\ref{tab:r7sweep}
to within sampling noise at every split (e.g.\ $0.650$ vs.\ $0.648$ at $a=10$,
$0.709$ vs.\ $0.709$ at $a=20$) --- the two independently-computed views of the
same model agree exactly, as they must.

\subsection{The geometry of the read-out matrix}
\label{sec:res-wout}

If the model relied on a fixed node-index shortcut (hypothesis~3), we would
expect a target-specific classifier vector $w_j$ to look unusually different
from its neighbours, or the direct read-in-to-read-out skip connection
$E_{\mathrm{in}}W_{\mathrm{out}}^\top$ (bypassing both transformer layers
entirely) to already carry most of the signal. Neither holds
(Table~\ref{tab:r7wout}): the $w_j$ row norms are uniform across all $40$
targets (standard deviation $\approx4\%$ of the mean, every seed), the
pairwise cosine similarity between different $w_j$'s is small
($\approx0.06$), and the skip-connection alignment is small everywhere
(diagonal $\approx0.2$--$0.3$, off-diagonal $\approx0$) --- two orders of
magnitude below the $h_i^\top w_j$ values of $10$--$35$ that actually decide a
prediction (\S\ref{sec:res-readout}). Whatever produces the large, split-dependent
logit swings is built by the two attention/feed-forward layers, not read
directly off which node is whose neighbour.

\begin{table}[H]\centering\small
\begin{tabular}{lcccc}
\toprule
 & $\|w_j\|$ (mean, std) & mean $\cos(w_j,w_k)$ & \multicolumn{2}{c}{$E_{\mathrm{in}}W_{\mathrm{out}}^\top$}\\
 & & ($j\neq k$) & diagonal & off-diagonal\\
\midrule
seed 1000 & $5.18\pm0.23$ & $+0.065$ & $+0.31$ & $-0.03\pm0.17$\\
seed 2000 & $5.21\pm0.19$ & $+0.063$ & $+0.25$ & $-0.06\pm0.19$\\
seed 3000 & $5.09\pm0.19$ & $+0.063$ & $+0.26$ & $-0.02\pm0.17$\\
seed 4000 & $5.12\pm0.20$ & $+0.061$ & $+0.24$ & $+0.03\pm0.16$\\
\bottomrule\end{tabular}
\caption{\emph{Model:} the four disjoint-paths-trained checkpoints of
Table~\ref{tab:r7sweep} (no test graphs --- these are static properties of the
trained weights). \emph{Metrics:} the read-out row norms $\|w_j\|$ (mean $\pm$
std.\ over the $40$ target nodes); the mean pairwise cosine similarity between
different rows of $W_{\mathrm{out}}$; and the read-in/read-out skip-connection
alignment $E_{\mathrm{in}}W_{\mathrm{out}}^\top$ (row $k$ of the read-in weight
against row $j$ of the read-out weight), split into its diagonal ($k=j$) and
off-diagonal mean $\pm$ std.}
\label{tab:r7wout}
\end{table}

\subsection{The same read-out geometry check on the ER-trained model}
\label{sec:res-wout-er}

The same static check on the ER-trained checkpoints of \S\ref{sec:res-sweep-er} gives the
same answer (Table~\ref{tab:r7wouter}): row norms uniform across targets, small pairwise
cosine similarity, and a skip-connection alignment an order of magnitude below the logit
values that actually decide a prediction. A different training distribution, entirely
different split-sweep behaviour (\S\ref{sec:res-sweep-er}), and the same refutation of a
fixed node-index shortcut --- whatever each model does with the split, neither reads it off
$W_{\mathrm{out}}$ or $W_{\mathrm{in}}$ directly.

\begin{table}[H]\centering\small
\begin{tabular}{lcccc}
\toprule
 & $\|w_j\|$ (mean, std) & mean $\cos(w_j,w_k)$ & \multicolumn{2}{c}{$E_{\mathrm{in}}W_{\mathrm{out}}^\top$}\\
 & & ($j\neq k$) & diagonal & off-diagonal\\
\midrule
seed 1000 & $5.68\pm0.16$ & $+0.127$ & $-0.52$ & $-0.16\pm0.33$\\
seed 2000 & $5.64\pm0.16$ & $+0.121$ & $-0.40$ & $-0.13\pm0.35$\\
seed 3000 & $5.61\pm0.17$ & $+0.125$ & $-0.32$ & $-0.10\pm0.33$\\
seed 4000 & $5.67\pm0.21$ & $+0.126$ & $-0.63$ & $-0.19\pm0.35$\\
\bottomrule\end{tabular}
\caption{\emph{Model:} the four ER-trained checkpoints of Table~\ref{tab:r7sweeper} (no test
graphs). \emph{Metrics as in Table~\ref{tab:r7wout}.}}
\label{tab:r7wouter}
\end{table}

\subsection{Attention: where the model actually looks}
\label{sec:res-attention}

\paragraph{Question.} The exact node-to-node contribution $C_{ik}=\|\partial
h_i^{(2)}/\partial h_k^{(0)}\|_F$ (\S\ref{sec:setup}) measures, for each query
node, how much its final representation actually moves with each other node's
read-in embedding. Does this quantity reach differently in a split the model
solves than in one it does not?

\paragraph{Result: the long component draws contribution from the short one
at every split tested, even the ones the model solves perfectly --- but far
more once the split balances.} For each long-component node we sum $C_{ik}$
over $k$ in its own (long) component and over $k$ in the short component, and
read the short share as a leak fraction. Figure~\ref{fig:r7attn} shows it is
never zero: even at $a=1$, where the model is exact on every graph, $5.8\%$ of
the long component's contribution mass already traces back to the single
short node. It grows from there, but \emph{not} smoothly --- it rises to
$11\%$ at $a=4$, $21$--$24\%$ at $a=7$--$8$, then spikes to $43\%$ right at
$a=10$ (immediately past where exact match collapses to zero), dips back to
$27\%$ at $a=14$, and climbs again to $31$--$36\%$ by $a=17$--$20$. Each of
these values is tight across the four seeds (a few percentage points of
spread at most), so the spike at $a=10$ is a reproducible feature of the
curve, not seed noise --- the mechanism does not de-confine itself gradually
and monotonically as the split balances; it is already leaking a little at
the splits the model gets perfectly right, and the leak grows unevenly,
peaking sharply right where behaviour first breaks before settling into a
broad elevated plateau. \textbf{The qualitative reading survives even though
the shape is messier than a smooth curve would suggest: confinement to the
query's own component is a matter of degree everywhere, and it degrades
substantially, if unevenly, as the split balances.}

\begin{figure}[H]
    \centering
    \figorbox{runs/report7/report7_figs/r7_attention_leak.png}{0.72}
    \caption{\emph{Model/training as in Table~\ref{tab:r7sweep}}. \emph{Test
    set:} two-chain graphs at the splits shown, $40$ graphs per seed.
    \emph{Metric:} for each long-component node, the fraction of its exact
    node-to-node contribution mass $C_{ik}=\|\partial h_i^{(2)}/\partial
    h_k^{(0)}\|_F$ (\S\ref{sec:setup}, real automatic differentiation through
    the whole forward pass) that traces back to the short component rather
    than its own long component, mean over long-component nodes; small dots
    are individual seeds, the line is the mean over seeds. Red dashed line:
    capacity $3^L=9$.}
    \label{fig:r7attn}
\end{figure}

\subsection{The same leak fraction on the ER-trained model: smooth where the disjoint-paths one was not}
\label{sec:res-attention-er}

The identical measurement on the ER-trained checkpoints (Figure~\ref{fig:r7attner}) gives a
strikingly \emph{cleaner} curve than Figure~\ref{fig:r7attn}: the leak fraction rises
smoothly and monotonically, from exactly $0\%$ at $a=1$ (every one of the four seeds) through
$5\%$ ($a=4$), $11$--$18\%$ ($a=7$--$10$), to $21\%$ ($a=14$), $25\%$ ($a=17$) and $30\%$
($a=20$) --- no spike, no dip, a textbook gradual de-confinement as the split balances. The
disjoint-paths model's non-monotone spike at $a=10$ (\S\ref{sec:res-attention}) is therefore
not a generic property of how this leak fraction behaves near a capacity transition; it is
specific to whatever the disjoint-paths training additionally builds into the completion
mechanism. The ER model's overall contribution magnitude is also far smaller throughout
(consistent with its modest, never-saturating read-out logit, \S\ref{sec:res-sweep-er}): the
raw within-long contribution near the diagonal is $\approx4$--$4.5$ at both $a=4$ and $a=20$,
roughly an order of magnitude below the disjoint-paths model's $53$ at $a=4$
(\S\ref{sec:res-heatmaps-er}).

\begin{figure}[H]
    \centering
    \figorbox{runs/report7/report7_figs/r7_attention_leak_n40_er.png}{0.72}
    \caption{\emph{Model/metric as in Figure~\ref{fig:r7attn}} but trained on pure sparse
    Erd\H{o}s--R\'enyi ($\mathrm{ER}(40,0.05)$), four seeds, $n=40$. \emph{Test set:} two-chain
    graphs at the splits shown, $64$ graphs per seed.}
    \label{fig:r7attner}
\end{figure}

\subsection{Falsification test: does the model actually merge two unrelated large components?}
\label{sec:res-threeway}

\paragraph{Question.} If the model's mechanism in \S\ref{sec:res-sweep}--\ref{sec:res-attention}
were simply ``resolve whichever component is small, default everything else to
connected'' (hypothesis~2 in its crude form), then on a graph with \emph{three}
disjoint paths --- one small and two large, the two large ones genuinely
disconnected from each other --- the model should wrongly call the two large
components connected to each other (both read only as ``not the small one'').

\paragraph{Result: the crude shortcut is refuted, but so is a full recovery ---
the completion trick is specific to a two-component canvas.} It is not refuted
by a coin flip: at every small-component size tested and in every one of the
four seeds, the model correctly calls the two large components disconnected
($\mathrm{cut}=1.000$ throughout, Table~\ref{tab:r7threeway}) --- the ``which
component'' decision generalises cleanly to three components, not just two.
But the within-large-component reach does \emph{not} recover to the
near-perfect values \S\ref{sec:res-sweep} found for a single long path beside a
small stub: it sits at $0.71$--$0.86$, matching almost exactly the degraded
reach of a \emph{balanced two-way} split of a similar size (Table~\ref{tab:r7sweep}:
$0.71$ at the balanced $(20,20)$ split). Adding a small, cleanly-resolvable
third component does not rescue the two large components' internal
connectivity --- each is treated as if it were the ``long'' side of a balanced
two-way split, not as the sole \emph{other} component in the graph. The
completion trick found in \S\ref{sec:res-sweep} is specific to graphs with
exactly two components, one of them small; it is not a general ``not-small
therefore connected'' rule.

\begin{table}[H]\centering\small
\begin{tabular}{lccccccc}
\toprule
small & large$_1$ & large$_2$ & reach large$_1$ & reach large$_2$ & cut(S,L$_1$) & cut(S,L$_2$) & cut(L$_1$,L$_2$)\\
\midrule
1  & 19 & 20 & 0.737 & 0.710 & 1.000 & 1.000 & 1.000\\
2  & 19 & 19 & 0.737 & 0.737 & 1.000 & 1.000 & 1.000\\
4  & 18 & 18 & 0.765 & 0.765 & 1.000 & 1.000 & 1.000\\
7  & 16 & 17 & 0.825 & 0.794 & 1.000 & 1.000 & 1.000\\
8  & 16 & 16 & 0.824 & 0.825 & 1.000 & 1.000 & 1.000\\
10 & 15 & 15 & 0.855 & 0.855 & 1.000 & 1.000 & 1.000\\
\bottomrule\end{tabular}
\caption{\emph{Model/training as in Table~\ref{tab:r7sweep}}. \emph{Test set:}
three disjoint paths partitioning all $40$ nodes --- one small path (size shown)
and two large paths (sizes large$_1$, large$_2$), the two large paths
\emph{not} connected to each other; $300$ graphs per seed, independently
node-relabelled. \emph{Metrics} (mean over four seeds, tight in every cell):
pairwise reach inside each large path (target connected); pairwise accuracy
between the small path and large$_1$/large$_2$, and between the two large
paths (all three target disconnected --- ``cut(L$_1$,L$_2$)'' is the decisive
column).}
\label{tab:r7threeway}
\end{table}

\subsection{The same falsification test on the ER-trained model: the crude shortcut is \emph{not} refuted}
\label{sec:res-threeway-er}

\paragraph{Result: at $n=40$, an ER-trained model \emph{does} show the crude ``not the small
one'' merge the disjoint-paths model refuted.} Table~\ref{tab:r7threeway-er} is the opposite
pattern from Table~\ref{tab:r7threeway}. At a very small third component ($\mathrm{small}=1,2$)
the two large, mutually disconnected components are wrongly called connected $91$--$94\%$ of
the time ($\mathrm{cut}(\mathrm{L}_1,\mathrm{L}_2)=0.094,0.062$); the error shrinks steadily
as the small component grows, reaching $0.82$ by $\mathrm{small}=10$ --- the exact shape
hypothesis~2's crude form predicts, and the same direction already seen for the
disjoint-paths model but only at the larger, sparser $n=64$ canvas (\S\ref{sec:res-n64}).
This is not a separate failure: it is the same over-connection already visible in the
two-way sweep of \S\ref{sec:res-sweep-er}, where cut degraded steadily from $1.000$ to
$0.149$ as the split balanced. Whether the ``which component'' skill survives a third,
disconnected piece is therefore not tied to which training distribution was used --- it
tracks whether \emph{that particular} checkpoint has an over-connection tendency at all, which
this report finds in the ER-trained model at $n=40$ and in the disjoint-paths model only at
the larger $n=64$ canvas, but not in the disjoint-paths model at $n=40$ or the ER model at
$n=64$.

\begin{table}[H]\centering\small
\begin{tabular}{lccccccc}
\toprule
small & large$_1$ & large$_2$ & reach large$_1$ & reach large$_2$ & cut(S,L$_1$) & cut(S,L$_2$) & cut(L$_1$,L$_2$)\\
\midrule
1  & 19 & 20 & 0.979 & 0.975 & 1.000 & 1.000 & 0.094\\
2  & 19 & 19 & 0.984 & 0.987 & 1.000 & 1.000 & 0.062\\
4  & 18 & 18 & 0.998 & 0.997 & 1.000 & 1.000 & 0.014\\
7  & 16 & 17 & 0.955 & 0.930 & 0.758 & 0.775 & 0.371\\
8  & 16 & 16 & 0.901 & 0.900 & 0.853 & 0.854 & 0.619\\
10 & 15 & 15 & 0.881 & 0.880 & 0.887 & 0.890 & 0.820\\
\bottomrule\end{tabular}
\caption{\emph{Model/test/metrics as in Table~\ref{tab:r7threeway}} but trained on pure
sparse Erd\H{o}s--R\'enyi ($\mathrm{ER}(40,0.05)$), four seeds, $n=40$.}
\label{tab:r7threeway-er}
\end{table}

\subsection{Attention and read-out weights, visualised}
\label{sec:res-heatmaps}

\paragraph{The attention pattern itself changes shape between a solved and a
failed split.} Figure~\ref{fig:r7attnscores} shows the real attention scores
$S=QK^\top/\sqrt{d_h}$ and the resulting normalized-ReLU $\alpha$, per layer,
for one example graph at $a=4$ and one at $a=20$. Layer $0$ is essentially
identical in both: a narrow band along the diagonal (each node attending to
its immediate path neighbours), broken cleanly at the component boundary ---
a purely local, structure-driven pattern that does not depend on the split.
Layer $1$ is where the two conditions diverge. The exact contribution
$C_{ik}$ (Figure~\ref{fig:r7rollout}, \S\ref{sec:setup}) is concentrated in a
band around the diagonal that \emph{decays} with distance rather than forming
a flat, uniform block --- at $a=4$, a pair three-or-fewer hops apart in the
long component averages roughly twice the contribution of a pair $15$ or more
hops apart ($53$ vs.\ $25$ in the units of Figure~\ref{fig:r7rollout}), and
the same ratio holds, at a much smaller overall scale, at $a=20$ ($8.5$ vs.\
$5.8$). What changes with the split is where that band is cleanly \emph{cut}:
at $a=20$ it breaks sharply at the component boundary into two separate
diagonal blocks, matching the leak-fraction curve of Figure~\ref{fig:r7attn};
at $a=4$ no such clean break is visible, because the short component is too
small (four nodes) for a break to show up as a distinct block against the
long one. One structural feature holds at \emph{every} split: the far
endpoint of each path --- the lone degree-$1$ node, structurally distinct
from every other node on the path --- is an unusually large \emph{source} of
contribution to the rest of its component (column sums $3$--$4\times$ the
per-node average at $a=4$ and $a=20$ alike) without being an unusually large
\emph{receiver}, consistent with the endpoint broadcasting a ``this path
closed'' signal outward rather than absorbing one. The overall scale itself
also tracks the split: raw contribution values are roughly $6$--$8\times$
larger at $a=4$ than at $a=20$ (colour scales in Figure~\ref{fig:r7rollout}
are not shared, on purpose, to keep each panel legible), the same direction
as the saturated-vs-faint raw logit magnitudes already seen in
\S\ref{sec:res-sweep}. The row-mass panels (right column,
Figure~\ref{fig:r7rollout}; unaffected by any of the above, since row-mass is
a property of $\alpha$ alone) make the attention-pattern shape concrete: at
$a=4$, layer-$1$ row-mass rises once past the component boundary and stays on
a flat, high plateau for the whole long path; at $a=20$, it rises after
\emph{each} boundary but dips again in the middle of each half, a double-hump
rather than a plateau. Per-node $Q/K/V$ vector norms (not shown in the main
figures; computed for every split) are essentially \emph{identical} in shape
between $a=4$ and $a=20$ --- the difference is not in how strongly a node's
query/key/value is scaled, only in which other nodes' keys and values it
lines up with.

\begin{figure}[H]
    \centering
    \figorbox{runs/report7/report7_figs/r7_heatmap_attention_scores.png}{0.95}
    \caption{\emph{Model:} one seed of the checkpoints of Table~\ref{tab:r7sweep}
    (base standard transformer, linear read-out, $L=2$, single head,
    $d_{\mathrm{model}}=512$, disjoint-paths-trained, $n=40$). \emph{Test set:}
    one representative two-chain graph at $a=4$ and one at $a=20$, node order
    fixed to base graph position (short path first, then long) for readability.
    \emph{Metric:} the real attention score matrix $S=QK^\top/\sqrt{d_h}$ and
    the resulting normalized-ReLU attention $\alpha=\mathrm{ReLU}(S)/n$, layer
    $0$ (top) and layer $1$ (bottom), averaged over $80$ independent random
    relabellings of the same underlying graph then mapped back to base node
    order (so the figure reflects the model's stable behaviour on this graph
    shape, not one noisy instance).}
    \label{fig:r7attnscores}
\end{figure}

\begin{figure}[H]
    \centering
    \figorbox{runs/report7/report7_figs/r7_heatmap_rollout_contrib.png}{0.95}
    \caption{\emph{Model/test as in Figure~\ref{fig:r7attnscores}}, $40$
    independent relabellings averaged for the top row. \emph{Metric:} the
    exact node-to-node contribution $C_{ik}=\|\partial h_i^{(2)}/\partial
    h_k^{(0)}\|_F$ (top row, \S\ref{sec:setup}; real automatic differentiation
    through $V$, $W_O$, the residual connections and the MLP, not an
    attention-only approximation) and the layer-$1$ message contribution
    $\alpha_{ij}\|W_O v_j\|$ (bottom row) --- how much node $j$'s value, after
    the output projection, actually adds to node $i$'s residual stream at
    that one layer --- for $a=4$ (left) and $a=20$ (middle); and the row-mass
    $\sum_j\alpha_{ij}$ (right column) along the chain, layer $0$ vs.\ layer
    $1$, for both splits, dotted line marks the component boundary. Colour
    scales are independent per panel (not shared across $a=4$ and $a=20$).}
    \label{fig:r7rollout}
\end{figure}

\paragraph{The read-out and read-in weight matrices show no exploitable
structure (further evidence against hypothesis~3).} Figure~\ref{fig:r7wout}
renders what Table~\ref{tab:r7wout} already summarised in numbers: $W_{\mathrm{out}}$
and $W_{\mathrm{in}}$ look like unstructured noise at the level of individual
entries, the cosine-similarity matrices between different target/label vectors
are uniformly small off the diagonal, and the skip-connection alignment
$E_{\mathrm{in}}W_{\mathrm{out}}^\top$ has no diagonal or banded pattern beyond
a few scattered entries --- there is no visible ``special'' node index, and no
sign that the model precomputes an index-based shortcut anywhere in these two
matrices. The singular-value spectra of every weight matrix in the trunk
(Figure~\ref{fig:r7sv}) decay smoothly with no isolated dominant direction,
consistent with a distributed computation rather than a single low-rank
shortcut circuit.

\begin{figure}[H]
    \centering
    \figorbox{runs/report7/report7_figs/r7_heatmap_weight_geometry.png}{0.95}
    \caption{\emph{Model:} the four checkpoints of Table~\ref{tab:r7wout} (static
    weights, no test graphs). \emph{Metric:} $W_{\mathrm{out}}$ and
    $E_{\mathrm{in}}=W_{\mathrm{in}}^\top$ (raw entries); the skip-connection
    alignment $E_{\mathrm{in}}W_{\mathrm{out}}^\top$; the cosine-similarity
    matrices $\cos(w_j,w_k)$ and $\cos(e_k,e_l)$; and the row norms
    $\|w_j\|,\|e_k\|$ by target/label index (one representative seed shown).}
    \label{fig:r7wout}
\end{figure}

\begin{figure}[H]
    \centering
    \figorbox{runs/report7/report7_figs/r7_heatmap_qkvo_weights.png}{0.75}
    \figorbox{runs/report7/report7_figs/r7_heatmap_singular_values.png}{0.62}
    \caption{\emph{Model as in Figure~\ref{fig:r7wout}.} \emph{Top:} the raw
    attention projection matrices $W_Q,W_K,W_V,W_O$ for both layers
    ($d_{\mathrm{model}}\times d_{\mathrm{model}}$). \emph{Bottom:} singular
    value spectra of every trunk weight matrix, normalised to the largest
    singular value (one representative seed).}
    \label{fig:r7sv}
\end{figure}

\subsection{The same heatmaps on the ER-trained model: a cleaner cut, no broadcasting endpoint}
\label{sec:res-heatmaps-er}

Figure~\ref{fig:r7rollout-er} repeats the exact-contribution heatmap on the ER-trained
checkpoints. Two differences from the disjoint-paths model stand out. First, the
near-diagonal band is already visibly \emph{cut} at the component boundary at $a=4$, not only
at $a=20$: a distinct bright block covers the four short-component nodes, separated from the
long component's own diagonal band, where the disjoint-paths model showed no such
separation until the short side grew much smaller a fraction of the canvas. Second, the
far-endpoint-as-source effect found throughout the disjoint-paths model (\S\ref{sec:res-heatmaps})
is \emph{absent} here: the two path endpoints' contribution-column sums are within a few
percent of the per-node average at both $a=4$ and $a=20$ (e.g.\ $105$ and $70$ against a mean
of $105$ at $a=4$), not the $3$--$4\times$ elevation seen in the disjoint-paths model. The raw
magnitude is also far smaller throughout, consistent with \S\ref{sec:res-attention-er}'s
reading of a model that never saturates its read-out logit: near-diagonal contribution is
$4.4$ at $a=4$ and $4.4$ at $a=20$ (compare $53$ and $8.5$ for the disjoint-paths model), a
roughly constant scale rather than the $6$--$8\times$ swing with split seen there. Both
differences point the same way: the disjoint-paths training builds a specific, split-shape-
sensitive signal (a broadcasting endpoint, a scale that tracks whether the short side is
resolvable) that an ER-trained model, which never saw a path endpoint mean anything, simply
does not acquire --- it falls back on the same modest, distance-decaying, cleanly-cut local
pattern at every split.

\begin{figure}[H]
    \centering
    \figorbox{runs/report7/report7_figs/r7_heatmap_rollout_contrib_n40_er.png}{0.95}
    \caption{\emph{Model/metric as in Figure~\ref{fig:r7rollout}} but trained on pure sparse
    Erd\H{o}s--R\'enyi ($\mathrm{ER}(40,0.05)$), $64$ independent relabellings averaged for
    the top row, $n=40$.}
    \label{fig:r7rollout-er}
\end{figure}

\subsection{Localising the completion signal: layer ablations}
\label{sec:res-ablation}

\paragraph{Question.} Which part of the trunk is responsible for the extra,
split-specific completion signal --- and is it the same machinery that does
ordinary distance-bounded reach, or a separate add-on?

\paragraph{Setup.} Starting from an otherwise normal forward pass, we zero one
branch of the residual stream at a time --- the attention output or the
feed-forward output, at layer $0$ or layer $1$ --- and re-run the full split
sweep; a fifth condition bypasses both transformer blocks entirely (read-out
applied directly to the read-in embedding, no cross-node mixing at all). If
ablating a component collapses the $a\le7$ behaviour toward the ablated
baseline everywhere, that component is doing the work; if the ablated model
is essentially unaffected, it is not.

\paragraph{Result: reach and cut depend on different parts of the trunk, and
the completion signal is not a separate circuit --- it needs the same
attention layers as ordinary reach.} Figure~\ref{fig:r7ablation} and
Table~\ref{tab:r7ablation} show a clean split. \textbf{Either} attention layer
alone is necessary for reach: zeroing layer $0$'s attention collapses reach to
$0.06$--$0.11$ at \emph{every} split, including $a\le7$; zeroing layer $1$'s
collapses it to $0.15$--$0.29$ --- both far below even the ordinary
beyond-capacity floor, and both close to the no-mixing-at-all bypass baseline
($0.32$--$0.34$, flat in $a$). The $a\le7$ completion signal does not survive
losing either attention layer any better than the ordinary $a\ge11$ regime
does: it is built by the \emph{same} two-layer attention machinery, not a
separate mechanism riding on top of it. The feed-forward branches tell a
different, sharper story. Zeroing layer $1$'s feed-forward leaves reach almost
untouched ($0.87$--$0.99$, barely below baseline at every split) but collapses
\emph{cut} from $\approx1.0$ to $0.26$--$0.62$ once the split is no longer
tiny --- \textbf{cut and reach are read out through different final-layer
computations}: cut depends heavily on layer $1$'s feed-forward step, reach
barely does. Layer $0$'s feed-forward instead matters more for reach in the
completion regime specifically (baseline $1.0\to0.47$ at $a=7$ under
ablation, a much larger relative hit than the $0.71\to0.59$ it takes at
$a=20$) while cut is only mildly affected throughout.

\begin{table}[H]\centering\small
\begin{tabular}{lcccccccc}
\toprule
condition & $a{=}1$ & $a{=}4$ & $a{=}7$ & $a{=}8$ & $a{=}10$ & $a{=}14$ & $a{=}17$ & $a{=}20$\\
\midrule
\multicolumn{9}{l}{\emph{reach (long component)}}\\
baseline    & 1.000 & 1.000 & 1.000 & 0.993 & 0.651 & 0.581 & 0.638 & 0.709\\
zero attn 0 & 0.058 & 0.063 & 0.068 & 0.070 & 0.074 & 0.084 & 0.094 & 0.107\\
zero attn 1 & 0.154 & 0.168 & 0.182 & 0.187 & 0.199 & 0.226 & 0.253 & 0.288\\
zero mlp 0  & 0.808 & 0.635 & 0.469 & 0.428 & 0.438 & 0.490 & 0.535 & 0.593\\
zero mlp 1  & 0.992 & 0.934 & 0.925 & 0.916 & 0.912 & 0.899 & 0.885 & 0.869\\
bypass      & 0.322 & 0.325 & 0.326 & 0.328 & 0.326 & 0.332 & 0.332 & 0.336\\
\midrule
\multicolumn{9}{l}{\emph{cut (target: disconnected)}}\\
baseline    & 1.000 & 1.000 & 0.998 & 0.926 & 0.999 & 1.000 & 1.000 & 1.000\\
zero attn 0 & 0.995 & 0.995 & 0.994 & 0.994 & 0.993 & 0.993 & 0.993 & 0.993\\
zero attn 1 & 0.986 & 0.997 & 0.996 & 0.996 & 0.995 & 0.995 & 0.993 & 0.994\\
zero mlp 0  & 0.969 & 0.887 & 0.898 & 0.945 & 0.942 & 0.917 & 0.914 & 0.912\\
zero mlp 1  & 0.939 & 0.599 & 0.617 & 0.618 & 0.534 & 0.276 & 0.262 & 0.256\\
bypass      & 0.660 & 0.692 & 0.689 & 0.693 & 0.688 & 0.691 & 0.692 & 0.691\\
\bottomrule\end{tabular}
\caption{\emph{Model/training as in Table~\ref{tab:r7sweep}}, four seeds.
\emph{Test set:} two-chain graphs at the splits shown, $150$ graphs per seed,
independently node-relabelled. \emph{Metric:} reach (long component, top block)
and cut (bottom block), mean over seeds, under six forward-pass conditions ---
unmodified (baseline); the attention or feed-forward branch zeroed at layer
$0$ or $1$; and both transformer blocks skipped entirely (bypass, read-out
applied directly to the read-in embedding).}
\label{tab:r7ablation}
\end{table}

\begin{figure}[H]
    \centering
    \figorbox{runs/report7/report7_figs/r7_ablation_reach_cut.png}{0.95}
    \caption{Table~\ref{tab:r7ablation} as curves. \emph{Model/training/test as
    in Table~\ref{tab:r7ablation}}, mean over four seeds. Red dashed line:
    capacity $3^L=9$.}
    \label{fig:r7ablation}
\end{figure}

\subsection{The same ablations on the ER-trained model: the double-edged mechanism, already at $n=40$}
\label{sec:res-ablation-er}

\paragraph{Result: zeroing either attention layer collapses reach but \emph{fixes} cut ---
the over-connection is built by the same machinery as the disjoint-paths model's completion
signal.} Table~\ref{tab:r7ablationer} repeats the ablation on the ER-trained checkpoints.
Reach depends on both attention layers just as it does for the disjoint-paths model (zeroing
layer $0$ collapses it to $0.04$--$0.07$, zeroing layer $1$ to $0.13$--$0.25$). Cut tells the
sharper story: zeroing \emph{either} attention layer pushes cut back up to $\ge0.996$ at every
split --- fully repairing the degraded baseline ($0.15$--$0.53$ once the split is no longer
tiny) --- while zeroing layer~$1$'s feed-forward collapses cut further still (down to $0.13$
at $a=20$), the same direction as the disjoint-paths model. This is exactly the
double-edged pattern this report finds for the disjoint-paths model only at the larger
$n=64$ canvas (\S\ref{sec:res-n64}): the identical two-layer attention computation that lets
the model confirm a connection is also what erodes its ability to keep two components
apart, and switching it off trades one failure for the other rather than fixing either
independently. Finding it here, in a model trained on a completely different distribution at
the \emph{smaller} $n=40$ canvas, says this is not a $64$-node-canvas-specific quirk: it is a
general property of how this two-layer attention mechanism trades reach against cut whenever
a checkpoint leans toward over-connection, whatever produced that lean.

\begin{table}[H]\centering\small
\begin{tabular}{lcccccccc}
\toprule
condition & $a{=}1$ & $a{=}4$ & $a{=}7$ & $a{=}8$ & $a{=}10$ & $a{=}14$ & $a{=}17$ & $a{=}20$\\
\midrule
\multicolumn{9}{l}{\emph{reach (long component)}}\\
baseline    & 0.965 & 0.999 & 0.977 & 0.946 & 0.902 & 0.925 & 0.942 & 0.960\\
zero attn 0 & 0.036 & 0.042 & 0.044 & 0.046 & 0.049 & 0.055 & 0.062 & 0.070\\
zero attn 1 & 0.131 & 0.141 & 0.154 & 0.159 & 0.169 & 0.194 & 0.218 & 0.251\\
zero mlp 0  & 0.440 & 0.511 & 0.456 & 0.458 & 0.465 & 0.500 & 0.531 & 0.577\\
zero mlp 1  & 0.994 & 0.981 & 0.973 & 0.972 & 0.968 & 0.959 & 0.949 & 0.932\\
bypass      & 0.231 & 0.231 & 0.232 & 0.232 & 0.232 & 0.231 & 0.231 & 0.233\\
\midrule
\multicolumn{9}{l}{\emph{cut (target: disconnected)}}\\
baseline    & 1.000 & 0.999 & 0.534 & 0.558 & 0.520 & 0.290 & 0.192 & 0.152\\
zero attn 0 & 1.000 & 0.999 & 0.998 & 0.997 & 0.997 & 0.996 & 0.996 & 0.996\\
zero attn 1 & 0.996 & 1.000 & 1.000 & 1.000 & 1.000 & 1.000 & 1.000 & 1.000\\
zero mlp 0  & 1.000 & 0.772 & 0.884 & 0.882 & 0.884 & 0.876 & 0.875 & 0.869\\
zero mlp 1  & 0.948 & 0.507 & 0.338 & 0.266 & 0.174 & 0.143 & 0.130 & 0.129\\
bypass      & 0.691 & 0.753 & 0.760 & 0.760 & 0.763 & 0.763 & 0.763 & 0.763\\
\bottomrule\end{tabular}
\caption{\emph{Model/test/metric as in Table~\ref{tab:r7ablation}} but trained on pure sparse
Erd\H{o}s--R\'enyi ($\mathrm{ER}(40,0.05)$), four seeds, $n=40$.}
\label{tab:r7ablationer}
\end{table}

\begin{figure}[H]
    \centering
    \figorbox{runs/report7/report7_figs/r7_ablation_reach_cut_n40_er.png}{0.95}
    \caption{Table~\ref{tab:r7ablationer} as curves. \emph{Model/training/test as in
    Table~\ref{tab:r7ablationer}}, mean over four seeds. Red dashed line: capacity $3^L=9$.}
    \label{fig:r7ablationer}
\end{figure}

\subsection{Generalising to $n=64$: path\_union and ER}
\label{sec:res-n64}

\paragraph{Question.} Every result so far is on one canvas size, $n=40$. Does the
same completion signal appear at a larger, sparser canvas, and does it depend
on having seen disjoint-path structure in training at all? We repeat the full
Tier~1/2 audit on the $n=64$ checkpoints from Report~VI's Thread~A --- already
trained, no new training needed --- for two training distributions: the same
disjoint-paths family as everywhere else in this report, and pure sparse
Erd\H{o}s--R\'enyi (which never contains a multi-path or two-chain graph).

\paragraph{Result: ER shows no completion signal at all, and tracks the
distance oracle almost exactly at every split.} Figure~\ref{fig:r7n64er} is the
cleanest single result in this report. Trained purely on ER, the model shows
\emph{no} special behaviour at small $a$: reach climbs smoothly from $0.21$ to
$0.39$ as $a$ goes from $1$ to $32$, sitting within a few points of the
distance-$\le9$ oracle rate at \emph{every} split, with no cliff, no floor, no
recovery --- exactly the profile of plain distance-bounded matrix powering,
never departing from it. Cut, meanwhile, stays at $\ge0.997$ across the entire
sweep, including the fully balanced split. The completion signal found
throughout this report is not an architectural inevitability of the trunk ---
it is something the disjoint-paths training distribution specifically teaches;
a model that never saw that structure does not acquire it, at any split.

\begin{figure}[H]
    \centering
    \figorbox{runs/report7/report7_figs/r7_sweep_and_logit_n64_er.png}{0.92}
    \caption{\emph{Model:} base standard transformer, linear read-out ($L=2$,
    single head, $d_{\mathrm{model}}=512$), trained on pure sparse
    Erd\H{o}s--R\'enyi ($p=0.03$, mean degree $\approx1.9$; never any path
    structure), four seeds, $n=64$ (checkpoints from Report~VI Thread~A).
    \emph{Test set/metrics as in Figure~\ref{fig:r7sweep}}, full split sweep
    $a=1,\dots,32$.}
    \label{fig:r7n64er}
\end{figure}

\paragraph{Result: path\_union at $n=64$ reproduces the completion signal for
reach, but the failure mode beyond it inverts --- cut collapses instead of
staying robust.} Figure~\ref{fig:r7n64pu} shows the disjoint-paths-trained
$n=64$ model does still single out small $a$ (exact match is high and reach
saturates near $1.0$ for $a\le4$, tracking the same qualitative shape as
$n=40$), but two things differ from $n=40$. First, reach \emph{never really
collapses}: it decays smoothly from $\approx1.0$ at $a=1$ to $\approx0.82$ at
$a=32$, always well above the distance oracle, with none of the sharp
$n=40$-style cliff --- on this larger, sparser canvas the model leans toward
\emph{over}-connecting rather than under-connecting once capacity is
exceeded. Second, and directly following from that: cut degrades steadily and
substantially, from $\approx0.9$ at small $a$ down to $\approx0.24$ by
$a=20$--$32$ --- the opposite of $n=40$, where cut stayed at $\approx1.0$
throughout. The raw read-out logit for cut (bottom panel) makes this precise:
it actually \emph{crosses zero} and turns mildly positive past $a\approx13$,
not merely fainter as at $n=40$ --- the model is not just less confident that
distant pairs are disconnected, for many of them it now positively believes
they are connected. Exact match at very small $a$ is also visibly noisier than
the tight $n=40$ result (per-seed at $a=4$: $0.43,\,0.55,\,0.79,\,0.00$),
consistent with Report~VI's own note that the $n=64$ path\_union checkpoints
are a noisier, more seed-dependent read than $n=40$'s tight clean model.

\begin{figure}[H]
    \centering
    \figorbox{runs/report7/report7_figs/r7_sweep_and_logit_n64_pathunion.png}{0.92}
    \caption{\emph{Model/metrics as in Figure~\ref{fig:r7n64er}} but trained on
    the same disjoint-paths distribution as every other result in this report,
    four seeds, $n=64$ (checkpoints from Report~VI Thread~A). Full split sweep
    $a=1,\dots,32$.}
    \label{fig:r7n64pu}
\end{figure}

\paragraph{The three-way falsification test flips too.} Repeating
\S\ref{sec:res-threeway} at $n=64$ (Table~\ref{tab:r7n64threeway}) reproduces
the ER-vs-path\_union contrast cleanly. The ER-trained model keeps
$\mathrm{cut}(\mathrm{L}_1,\mathrm{L}_2)\ge0.998$ at every small-component
size, exactly as at $n=40$ --- the two large, unrelated components are
correctly kept apart regardless of training distribution or canvas size. The
path\_union-trained model does \emph{not}: at a very small third component
($\mathrm{small}=1$), it wrongly calls the two large ($31$- and $32$-node)
components connected $91\%$ of the time
($\mathrm{cut}(\mathrm{L}_1,\mathrm{L}_2)=0.093$) --- the crude ``not the small
one'' shortcut hypothesis~2 explicitly predicted and that $n=40$ refuted. The
error shrinks steadily as the small component grows ($0.09\to0.79$ from
$\mathrm{small}=1$ to $10$), the same direction as the two-way sweep's cut
degrading with more structure to keep separate. So hypothesis~2's crude form,
refuted at $n=40$, is \emph{not} refuted in general --- whether the ``which
component'' skill survives a third component depends on canvas size and
training distribution, tracking the same over-connection tendency visible in
Figure~\ref{fig:r7n64pu}.

\begin{table}[H]\centering\small
\begin{tabular}{lcccccc}
\toprule
 & \multicolumn{3}{c}{path\_union-trained} & \multicolumn{3}{c}{ER-trained}\\
\cmidrule(lr){2-4}\cmidrule(lr){5-7}
small & reach L$_1$/L$_2$ & cut(S,L) & cut(L$_1$,L$_2$) & reach L$_1$/L$_2$ & cut(S,L) & cut(L$_1$,L$_2$)\\
\midrule
1  & 0.926 & 0.999 & 0.093 & 0.396 & 1.000 & 0.998\\
2  & 0.707 & 1.000 & 0.369 & 0.404 & 1.000 & 0.998\\
4  & 0.756 & 0.999 & 0.322 & 0.417 & 1.000 & 0.998\\
7  & 0.503 & 0.997 & 0.612 & 0.431 & 1.000 & 1.000\\
8  & 0.427 & 0.997 & 0.717 & 0.438 & 0.999 & 0.999\\
10 & 0.394 & 0.998 & 0.785 & 0.453 & 1.000 & 1.000\\
\bottomrule\end{tabular}
\caption{\emph{Model/training/test as in Table~\ref{tab:r7threeway}}, $n=64$,
four seeds each. \emph{Metrics:} pairwise reach inside the two large components
(mean of L$_1$, L$_2$); cut between the small and either large component
(mean); cut between the two large components (target disconnected --- the
decisive column, contrast the two blocks).}
\label{tab:r7n64threeway}
\end{table}

\paragraph{Ablation confirms layer-$1$ attention is a double-edged mechanism at
$n=64$: the same computation that builds the completion signal also builds the
over-connection.} Figure~\ref{fig:r7n64ablation} repeats \S\ref{sec:res-ablation}
on the path\_union $n=64$ model. Zeroing layer-$1$ attention collapses reach
just as it does at $n=40$ ($0.05$--$0.10$ at every split) --- but it also
\emph{fixes} cut, pushing it to $\ge0.93$ at every split, well above the
degraded baseline. The same is true, more mildly, of zeroing layer-$0$
attention (cut steady at $\approx0.88$, versus baseline's fall to $0.24$).
Zeroing layer-$1$'s feed-forward again collapses cut specifically (down to
$0.15$ at $a=32$) while leaving reach largely intact, replicating
\S\ref{sec:res-ablation} exactly. Put together: the two-layer attention
mechanism that builds the beyond-capacity reach signal is the \emph{same}
mechanism responsible for the over-connection that erodes cut on this larger
canvas --- removing it trades one failure for the other rather than fixing
either independently.

\begin{figure}[H]
    \centering
    \figorbox{runs/report7/report7_figs/r7_ablation_reach_cut_n64_pathunion.png}{0.92}
    \caption{\emph{Model/training/test as in Table~\ref{tab:r7ablation}},
    $n=64$, path\_union-trained, four seeds. \emph{Metric as in
    Figure~\ref{fig:r7ablation}.}}
    \label{fig:r7n64ablation}
\end{figure}

\subsection{A third canvas size: $n=20$, where the capacity wall cannot bite}
\label{sec:res-n20}

\paragraph{Question.} \S\ref{sec:res-n64} generalised from $n=40$ to a larger, sparser
canvas. At $n=20$ the same architectural capacity ($3^L=9$) is close to the whole canvas: even
the most balanced split, $(10,10)$, has a within-component distance that never exceeds $9$.
So a genuine test of whether this report's mechanism is specific to $n=40$ should not expect
the same behavioural wall to reappear at $n=20$ --- it is a real, informative check
precisely \emph{because} it removes the capacity constraint that drives every other result in
this report, not a re-confirmation we expect to look identical.

\paragraph{Result: the disjoint-paths model solves every split, exactly as the capacity
argument predicts --- but the underlying leak fraction still rises with the split.}
Table~\ref{tab:r7n20sweep} confirms the expectation directly: the disjoint-paths-trained
model is exact on essentially every graph at every split from $a=1$ to the balanced
$a=10$ ($\ge0.98$ throughout, reach and cut both $1.000$) --- there is no wall to expose
because nothing tested ever asks it to confirm a pair beyond $9$ hops. The interesting result
is that the mechanistic signature from \S\ref{sec:res-attention} does not disappear along
with the behavioural wall: the leak fraction (Figure~\ref{fig:r7n20leak}, left) still rises
smoothly from $0.4\%$ at $a=1$ to $36\%$ at the balanced $a=10$, essentially the same range as
the $n=40$ result, even though behaviour here is perfect throughout. The completion mechanism
therefore is not simply a symptom of the capacity wall being hit; it is a graded property of
how the trunk handles a two-component graph that shows up whether or not the split ever
forces a visible mistake.

\begin{table}[H]\centering\small
\begin{tabular}{lccccc}
\toprule
split $(a,\,20{-}a)$ & exact & reach (long) & reach (short) & cut & pred.\ positive rate\\
\midrule
$(1,19)$  & 1.000 & 1.000 & --    & 1.000 & 0.905\\
$(4,16)$  & 1.000 & 1.000 & 1.000 & 1.000 & 0.680\\
$(7,13)$  & 0.987 & 1.000 & 1.000 & 1.000 & 0.545\\
$(8,12)$  & 0.984 & 1.000 & 1.000 & 1.000 & 0.520\\
$(10,10)$ & 1.000 & 1.000 & 1.000 & 1.000 & 0.500\\
\bottomrule\end{tabular}
\caption{\emph{Model:} base standard transformer, linear read-out ($L=2$, single head,
$d_{\mathrm{model}}=512$), trained on the same clean disjoint-paths distribution as
Table~\ref{tab:r7sweep}, four seeds, $n=20$. \emph{Test set/metrics as in
Table~\ref{tab:r7sweep}}, full sweep $a=1,\dots,10$ (selected columns shown).}
\label{tab:r7n20sweep}
\end{table}

\paragraph{Result: an ER-trained model at $n=20$ never solves the two-chain task well at any
split, for a reason that is not the capacity wall either.} Table~\ref{tab:r7n20sweeper} shows
a genuinely different picture from every other ER result in this report: exact match is low
and non-monotone across the whole sweep (peaking at $46\%$ around $a=3$, not tabulated, and
falling to $0$ by $a=6$--$9$ before a small recovery at $a=10$), and reach dips in the middle
of the sweep ($0.77$--$0.87$ at $a=3$--$9$) rather than at one end. Since nothing here exceeds
the capacity, this is not the distance wall reappearing; it reads instead as a general
out-of-distribution sensitivity to two-chain graphs at this canvas size, worst somewhere in
the middle of the split range rather than at the balanced point specifically. The leak
fraction (Figure~\ref{fig:r7n20leak}, right) still rises with the split, from $2\%$ at $a=1$
to $46\%$ at $a=10$, the same qualitative direction as every other condition in this report,
despite the noisier behavioural backdrop.

\begin{table}[H]\centering\small
\begin{tabular}{lccccc}
\toprule
split $(a,\,20{-}a)$ & exact & reach (long) & reach (short) & cut & pred.\ positive rate\\
\midrule
$(1,19)$  & 0.161 & 0.932 & --    & 0.995 & 0.848\\
$(4,16)$  & 0.176 & 0.836 & 0.988 & 0.990 & 0.584\\
$(7,13)$  & 0.000 & 0.842 & 0.922 & 0.904 & 0.519\\
$(8,12)$  & 0.000 & 0.852 & 0.912 & 0.887 & 0.513\\
$(10,10)$ & 0.060 & 0.886 & 0.887 & 0.868 & 0.515\\
\bottomrule\end{tabular}
\caption{\emph{Model:} base standard transformer, linear read-out ($L=2$, single head,
$d_{\mathrm{model}}=512$), trained on pure sparse Erd\H{o}s--R\'enyi ($\mathrm{ER}(20,0.08)$,
the paper's original density), four seeds, $n=20$. \emph{Test set/metrics as in
Table~\ref{tab:r7sweep}}, full sweep $a=1,\dots,10$ (selected columns shown).}
\label{tab:r7n20sweeper}
\end{table}

\begin{figure}[H]
    \centering
    \figorbox{runs/report7/report7_figs/r7_attention_leak_n20_pathunion.png}{0.48}
    \figorbox{runs/report7/report7_figs/r7_attention_leak_n20_er.png}{0.48}
    \caption{\emph{Model/metric as in Figure~\ref{fig:r7attn}}, $n=20$, four seeds, $64$ graphs
    per seed. \emph{Left:} disjoint-paths-trained (Table~\ref{tab:r7n20sweep}). \emph{Right:}
    ER-trained (Table~\ref{tab:r7n20sweeper}).}
    \label{fig:r7n20leak}
\end{figure}

\paragraph{The three-way falsification test is clean for both models here --- consistent with
over-connection being a beyond-capacity phenomenon.} Neither model wrongly merges the two
large, disconnected components at $n=20$: the disjoint-paths model is exact
($\mathrm{cut}(\mathrm{L}_1,\mathrm{L}_2)=1.000$ at every small-component size tested), and the
ER model is close behind ($0.95$--$1.00$). This is the cleanest single confirmation in this
report that the crude ``not the small one'' merge found for the ER model at $n=40$
(\S\ref{sec:res-threeway-er}) and for the disjoint-paths model at $n=64$
(\S\ref{sec:res-n64}) is a beyond-capacity failure mode specifically: remove the possibility
of exceeding capacity, and both training distributions keep three components straight.

\paragraph{Ablation: reach is more robust here, and cut depends on more than one component of
the trunk.} The double-edged pattern from \S\ref{sec:res-ablation}/\S\ref{sec:res-ablation-er}
is present but less clean-cut at $n=20$. For the disjoint-paths model, reach survives
attention ablation far better than at $n=40$ (zeroing layer~$0$ leaves reach at
$0.40$--$0.43$, zeroing layer~$1$ at $0.79$--$0.81$, against a baseline of $1.000$ ---
nowhere near the near-total collapse seen at $n=40$/$64$), while cut, which is perfect at
baseline, is knocked down by \emph{every} ablation tested ($0.38$--$0.85$), not specifically
by attention. For the ER model, reach again depends heavily on both attention layers
($0.11$--$0.17$ and $0.27$--$0.47$ under ablation), and zeroing either attention layer still
moves cut in the same repairing direction seen at $n=40$ (up to $0.99$ from a baseline of
$0.87$--$0.99$), though the effect is milder since the ER baseline here is already closer to
correct. Read together with \S\ref{sec:res-ablation-er}, the double-edged reach/cut trade-off
in layer-$1$ attention looks like a general feature of this architecture, not a $64$-node or
$40$-node-specific artefact --- it is just harder to see at $n=20$ because there is less room
for either failure mode to develop.

\subsection{Does the similarity read-out learn the same thing?}
\label{sec:res-similarity}

\emph{[Results pending --- the full battery of \S\ref{sec:res-sweep}--\ref{sec:res-ablation-er}
is being re-run, eval-only, on the existing similarity-read-out checkpoints (both training
distributions, $n=40$, four seeds each); no data collected yet, nothing below is written from
invented numbers.]}

\section{Verdict}
\label{sec:verdict}

We stated three candidate hypotheses in \S\ref{sec:hypotheses}. The mechanistic
evidence settles them as follows.

\paragraph{Hypothesis~3 (label/order shortcut) --- refuted.} Fixing node labels
to a canonical order instead of relabelling every graph independently changes
nothing (Table~\ref{tab:r7relabel}). The read-out matrix shows no
target-specific outlier structure, and the direct read-in-to-read-out skip
connection carries a signal two orders of magnitude smaller than what actually
decides a prediction (Table~\ref{tab:r7wout}) --- whatever the model does, it
is not reading node identity off a fixed index mapping.

\paragraph{Hypothesis~2 (short-component-resolved default) --- correct in a
narrower form than stated, and that form is itself canvas- and
training-distribution-dependent.} The model does not implement a generic ``not
in the small component therefore connected'' rule: at $n=40$, the three-way
falsification test shows it correctly separates two large, unrelated
components from a small one at $\mathrm{cut}=1.000$ throughout
(Table~\ref{tab:r7threeway}), which a crude default-connect rule would get
wrong. What survives at $n=40$ is narrower and, we think, more interesting: on
a graph with \emph{exactly two} components, once one of them is small enough
that the model can resolve it completely, the \emph{other} component gets an
extra, split-specific completion signal --- visible as a smoothly saturating
positive read-out logit (\S\ref{sec:res-readout}) and as attention that stays
cleanly confined to the true component (\S\ref{sec:res-attention}) --- that
goes beyond plain distance-bounded reach. That signal is switched off, not
merely weakened, once the graph has a third, separate component at $n=40$: the
two large components in \S\ref{sec:res-threeway} are each treated exactly like
the long side of a balanced two-way split, not like the sole beneficiary of a
completion rule. \textbf{This clean picture does not fully survive elsewhere in this report}, and not along the
line ``ER never shows it, disjoint-paths sometimes does'' that the $n=64$ result alone would
suggest. At $n=64$ (\S\ref{sec:res-n64}), the disjoint-paths-trained model wrongly merges two
large unrelated components when the third is very small
($\mathrm{cut}(\mathrm{L}_1,\mathrm{L}_2)=0.09$ at $\mathrm{small}=1$) while the ER-trained
model at the same canvas keeps cut $\ge0.998$ throughout. But at $n=40$ the roles reverse: the
disjoint-paths model is the clean one (Table~\ref{tab:r7threeway}), while an \emph{ER-trained}
model shows the crude merge just as badly ($\mathrm{cut}(\mathrm{L}_1,\mathrm{L}_2)=0.01$--$0.09$
at small third components, \S\ref{sec:res-threeway-er}). So the crude form of hypothesis~2 is
not tied to which training distribution was used; it tracks whether a given checkpoint has an
over-connection tendency \emph{at all} on that canvas --- something this report finds in the
disjoint-paths model only at $n=64$ and in the ER model only at $n=40$, matching each model's
own two-way cut degradation (\S\ref{sec:res-n64}, \S\ref{sec:res-sweep-er}) almost exactly.
\S\ref{sec:res-n20} pins down \emph{when} it can happen at all: at $n=20$, where no split can
exceed the architectural capacity, \emph{both} training distributions keep the three components
straight ($\mathrm{cut}(\mathrm{L}_1,\mathrm{L}_2)\ge0.95$ throughout). The crude merge is a
beyond-capacity failure mode, available to any checkpoint that has learned to over-connect once
distance exceeds $3^L$, not a property of one architecture or one training distribution.

\paragraph{Hypothesis~1 (genuine, distance-graded propagation) --- correct once
the completion signal is not active.} Beyond the two-component-with-a-small-side
regime, the model's aggregate behaviour is very hard to distinguish from plain
distance-bounded matrix powering: the overall predicted-positive rate for
$a\gtrsim11$ sits almost exactly on the rate a pure distance-$\le9$ oracle would
give on the same graphs (\S\ref{sec:res-sweep}), and the attention-leak curve
(Figure~\ref{fig:r7attn}) is graded and already non-zero at the splits the
model solves perfectly, rather than switching on and off at a threshold. The
``genuine propagation'' account is the right description of the model
\emph{outside} the specific regime where the completion shortcut fires --- it
is not a competing story to hypothesis~2, it is what the model falls back to
when hypothesis~2's narrow condition is not met.

\paragraph{Put together.} The Report~VI puzzle is not really about balance. It
is about whether a two-component graph contains one component small enough to
fully resolve. When it does, the model gets an extra completion signal for the
other component that reaches far beyond its ordinary distance-bounded capacity
(the $4{+}36$ case); when it does not --- whether because the split is
balanced, or because a third component is present at all --- that signal does
not fire, and the model's behaviour reverts to the same distance-bounded,
$3^L$-capped propagation documented since Report~III.

\paragraph{Where the signal lives.} The ablations of \S\ref{sec:res-ablation}
add a mechanistic layer to that picture rather than a competing one. Reach and
cut are computed by different parts of the trunk --- cut depends heavily on
layer $1$'s feed-forward step and barely on either attention layer's presence
\emph{once both are intact}, while reach collapses if \emph{either} attention
layer is removed, at every split, including the completion regime. In other
words, the $a\le7$ completion signal is not a separate sub-circuit bolted onto
an otherwise-ordinary matrix-powering computation: it needs the same
two-layer attention machinery that ordinary beyond-capacity reach needs, and
the attention-score heatmaps (\S\ref{sec:res-heatmaps}) show \emph{why} ---
the same layer-$1$ attention that forms a distance-decaying band cleanly cut
at the component boundary at large $a$ forms the same decaying band
\emph{without} that clean cut once the other side is small enough to
resolve, because a stub of only a few nodes is too small to register as a
separate block. It is the same mechanism doing more, not two mechanisms.

\paragraph{$n=64$ confirms the mechanism is data-shaped, not architectural, and
reveals its dark side.} The ER-vs-path\_union contrast at $n=64$
(\S\ref{sec:res-n64}) is the cleanest evidence in this report that the
completion signal is something specific training distributions teach: the
ER-trained model never acquires it, at any split, and tracks a plain
distance-bounded oracle almost exactly throughout. The path\_union-trained
model does acquire it, and at $n=64$ the ablations show it is a double-edged
version of the \emph{same} mechanism found at $n=40$ --- the identical
two-layer attention computation that lets a small-side graph confirm
connections far beyond capacity also, on a larger and sparser canvas, starts
declaring unrelated large components connected to each other. Zeroing
layer-$1$ attention removes both effects together (Figure~\ref{fig:r7n64ablation}):
reach collapses, but cut is fully restored. Reach and cut were already shown
to depend on different final-layer computations (cut on layer-$1$
feed-forward, largely independent of either attention layer's presence); what
$n=64$ adds is that the attention mechanism which \emph{builds} the useful
completion signal is the same one that, past a point, over-generalises it into
a harmful one.

\paragraph{The double-edged mechanism is general; one specific piece of it is not.} The
ER-trained model at $n=40$ (\S\ref{sec:res-ablation-er}) shows the identical double-edged
ablation signature found at $n=64$ --- zeroing either attention layer restores its already
badly-degraded cut to $\ge0.996$ while collapsing reach --- so this trade-off is not a
$64$-node-canvas-specific curiosity; it appears wherever a checkpoint has developed an
over-connection tendency, which \S\ref{sec:res-n20} shows requires exceeding capacity in the
first place. One piece of the $n=40$ disjoint-paths mechanism does \emph{not} generalise this
way: the far path-endpoint's role as an outsized \emph{source} of contribution
(\S\ref{sec:res-heatmaps}) is absent in the ER-trained model at the same canvas
(\S\ref{sec:res-heatmaps-er}), whose endpoints contribute at the same level as every other
node. A model that never saw a path end mean anything in training does not single that
position out --- so this specific broadcasting signal, unlike the general reach/cut trade-off
in layer-$1$ attention, is something the disjoint-paths distribution specifically teaches,
consistent with how this report reads the completion signal as a whole.

\section{Honest limits}
\label{sec:limits}

This report opens the trunk further than any before it, but the evidence
remains statistical and correlational, not a proof of what the network
computes. The read-out, attention, and ablation analyses show quantities and
sub-computations that move in lockstep with behaviour; they do not amount to
a formal account of \emph{why} training produces a completion signal specific
to a two-component canvas. A causal intervention --- transplanting activations
between a solved and a failing graph, or intervening on attention patterns
directly --- is left for future work. Two further follow-ups were not run
in this round and are left open: a retrain with a loss that reweights
within-component and cross-component pairs evenly, to separate a genuine
representational effect from a training-objective artefact (the confound
Report~VI's own outlook names); and the same audit repeated with the
similarity read-out, whose logit is a direct function of \emph{both}
endpoints' embeddings and may not show the same asymmetric, target-specific
completion signal at all. We have not checked whether the two-component
completion signal survives under the mixed training stream of earlier
reports --- the $n=64$ generalisation (\S\ref{sec:res-n64}) checks canvas size
and a from-scratch training distribution, not that stream specifically.

\end{document}
